\documentclass[10pt,a4paper]{amsart}
\usepackage[margin=19mm]{geometry}
\usepackage{amsmath,amssymb,mathtools}
\usepackage[scr=rsfs,cal=euler]{mathalfa}
\usepackage{xfrac}
\usepackage[unicode,hidelinks,bookmarks=false]{hyperref}
\newcommand{\ie}{i.e.\ }
\newcommand{\cf}{cf.\ }
\newcommand{\resp}{respectively\ }
\newcommand{\Subsection}{\subsection}
\providecommand{\nobreakdash}{\nobreak\hbox{-}}
\setlength{\emergencystretch}{2em}
\multlinegap0pt
\providecommand{\theHalgorithm}{\arabic{algorithm}}

% Some newcommands.
\newcommand{\norm}[1]{\left \lVert #1 \right \rVert}
\newcommand{\abs}[1]{\lvert #1\rvert}
\newcommand{\Abs}[1]{\left\lvert #1\right\rvert}

\providecommand{\bra}[1]{\langle #1|}
\providecommand{\ket}[1]{|#1 \rangle}
\providecommand{\braket}[2]{\langle #1|#2 \rangle}
\providecommand{\ketbra}[2]{|#1\rangle \langle #2|}
\providecommand{\expval}[1]{\langle #1 \rangle}
\newcommand{\dsA}{\mathbb{A}}
\newcommand{\dsB}{\mathbb{B}}
\newcommand{\dsC}{\mathbb{C}}
\newcommand{\dsD}{\mathbb{D}}
\newcommand{\dsE}{\mathbb{E}}
\newcommand{\dsF}{\mathbb{F}}
\newcommand{\dsG}{\mathbb{G}}
\newcommand{\dsH}{\mathbb{H}}
\newcommand{\dsI}{\mathbb{I}}
\newcommand{\dsJ}{\mathbb{J}}
\newcommand{\dsK}{\mathbb{K}}
\newcommand{\dsL}{\mathbb{L}}
\newcommand{\dsM}{\mathbb{M}}
\newcommand{\dsN}{\mathbb{N}}
\newcommand{\dsO}{\mathbb{O}}
\newcommand{\dsP}{\mathbb{P}}
\newcommand{\dsQ}{\mathbb{Q}}
\newcommand{\dsR}{\mathbb{R}}
\newcommand{\dsS}{\mathbb{S}}
\newcommand{\dsT}{\mathbb{T}}
\newcommand{\dsU}{\mathbb{U}}
\newcommand{\dsV}{\mathbb{V}}
\newcommand{\dsW}{\mathbb{W}}
\newcommand{\dsX}{\mathbb{X}}
\newcommand{\dsY}{\mathbb{Y}}
\newcommand{\dsZ}{\mathbb{Z}}
\newcommand{\scA}{\mathcal{A}}
\newcommand{\scB}{\mathcal{B}}
\newcommand{\scC}{\mathcal{C}}
\newcommand{\scD}{\mathcal{D}}
\newcommand{\scE}{\mathcal{E}}
\newcommand{\scF}{\mathcal{F}}
\newcommand{\scG}{\mathcal{G}}
\newcommand{\scH}{\mathcal{H}}
\newcommand{\scI}{\mathcal{I}}
\newcommand{\scJ}{\mathcal{J}}
\newcommand{\scK}{\mathcal{K}}
\newcommand{\scL}{\mathcal{L}}
\newcommand{\scM}{\mathcal{M}}
\newcommand{\scN}{\mathcal{N}}
\newcommand{\scO}{\mathcal{O}}
\newcommand{\scP}{\mathcal{P}}
\newcommand{\scQ}{\mathcal{Q}}
\newcommand{\scR}{\mathcal{R}}
\newcommand{\scS}{\mathcal{S}}
\newcommand{\scT}{\mathcal{T}}
\newcommand{\scU}{\mathcal{U}}
\newcommand{\scV}{\mathcal{V}}
\newcommand{\scW}{\mathcal{W}}
\newcommand{\scX}{\mathcal{X}}
\newcommand{\scY}{\mathcal{Y}}
\newcommand{\scZ}{\mathcal{Z}}
\newcommand{\Tr}{\operatorname{Tr}}
\newcommand{\ii}{\mathrm{i}}
\newcommand{\ee}{\mathrm{e}}
\newcommand{\dd}{\mathrm{d}}
\providecommand{\vr}{}\renewcommand{\vr}{\mathbf{r}}

\DeclareMathOperator{\Sym}{Sym}
\DeclareMathOperator{\End}{End}
\DeclareMathOperator{\TV}{TV}
\DeclareMathOperator{\KL}{KL}

% Author comments. Disable or remove for final submission as needed.
\newcommand{\QZ}[1]{{\textcolor{blue}{#1}}}
\newcommand{\JY}[1]{{\textcolor{orange}{#1}}}
\newcommand{\PL}[1]{{\textcolor{green}{#1}}}
\newcommand{\XC}[1]{{\textcolor{purple}{#1}}}

\theoremstyle{plain}
\newtheorem{theorem}{Theorem}
\newtheorem{proposition}[theorem]{Proposition}
\newtheorem{lemma}[theorem]{Lemma}
\newtheorem{corollary}[theorem]{Corollary}
\theoremstyle{definition}
\newtheorem{definition}[theorem]{Definition}
\newtheorem{assumption}[theorem]{Assumption}
\theoremstyle{remark}
\newtheorem{remark}[theorem]{Remark}
\title[Hierarchy of discriminative power and complexity in learning quantum ensembles]{Hierarchy of discriminative power and complexity in learning quantum ensembles: sufficiency of $k=N$ for full discriminative power (Theorem A.3)}
\author{Jian Yao, Pengtao Li, Xiaohui Chen, Quntao Zhuang}
\date{Source: arXiv:2601.22005v2 (CC BY 4.0)}
\begin{document}
\maketitle

\noindent\emph{Source and licence.} Hierarchy of discriminative power and complexity in learning quantum ensembles. Version arXiv:2601.22005v2 (CC BY 4.0), distributed under the Creative Commons Attribution 4.0 International licence, as recorded in the arXiv licence metadata for this exact version and shown on the abstract page saved with this item. The full licence notice, which carries the licence URL and the address of that abstract page, is the bundled file \texttt{licenses/arxiv-2601.22005-CC-BY-4.0.txt}.\par\smallskip

\noindent\textit{Editorial scope.} Counted unit: the article's Theorem A.3 (source0.tex lines 973--990, source label \texttt{thm: k=N full discriminative power}), ``Sufficiency: $k=N$ gives full discriminative power for ensembles with at most $N$ pure states'' -- the statement that two ensembles of at most $N$ distinct pure states whose $N$-th moment operators $M_N$ agree are equal up to a permutation of indices -- delivered with the article's complete original proof (source0.tex lines 992--1113, one \texttt{proof} environment of 122 lines, adjacent to the statement, with no gap and no deferral). No mathematics is written, repaired or re-derived: statement and proof are the article's own text line for line, in the article's own notation ($\scE_1,\scE_2$, $N_1,N_2$, $p_i,q_j$, $\ket{\psi_i},\ket{\phi_j}$, $M_t(\mathcal E)$, $D^{(k)}$, $\dsE$, $\Abs{\cdot}$, $\braket{\cdot}{\cdot}$), and no line of the counted statement or of the counted proof is omitted, substituted or re-flowed. Required context retained uncounted: the article's Lemma ``The hierarchy of discriminative power of MMD-$k$'' (source0.tex lines 379--381, source label \texttt{thm: hierarachy of MMD-k}), which the first line of the counted proof invokes; the article's Lemma ``Homogeneous polynomials and symmetric tensors'' (lines 828--845, source label \texttt{lemma: homogeneous and sym}), which supplies the polynomial-moment identity used in the proof; and the article's Lemma ``Hyperplane separation for distinct pure states'' (lines 936--947, source label \texttt{lemma: hyperplane}), which supplies the separating functional used in the proof. Editorial formatting only, in addition: an AMS \texttt{amsart} preamble that restates the article's own macro definitions (source0.tex lines 73--152, including \texttt{\textbackslash ket}, \texttt{\textbackslash braket}, \texttt{\textbackslash ketbra}, \texttt{\textbackslash dsE}, \texttt{\textbackslash scE}, \texttt{\textbackslash scP}, \texttt{\textbackslash Abs}, \texttt{\textbackslash Tr} and the article's \texttt{\textbackslash ds}- and \texttt{\textbackslash sc}-families) and the theorem environments the retained lines use; the article's own class (revtex4-2) and its shipped style files are neither shipped nor input, and the retained lines use no command that only those files provide; sequential numbering of the retained environments in order of appearance, so that the retained hierarchy lemma, polynomial lemma, hyperplane lemma and counted theorem are numbered 1, 2, 3 and 4, whereas the article prints them as Lemma A.1, Lemma A.2 (or its equivalents in the appendix) and Theorem A.3; the article's section-relative equation numbering is not reproduced, so displayed equations inside the retained spans are numbered consecutively instead, and every cross-reference inside the retained text resolves to the wrapper's numbering. No citation occurs inside any retained line, so the wrapper carries no bibliography and no bibliography entry. No asset-inclusion command, no drawing environment and no image asset occurs in this package; the article's nine figure files lie outside every retained span.

\section*{Retained context: the MMD-$k$ hierarchy and the two auxiliary lemmas}

\begin{lemma}[\textbf{The hierarchy of discriminative power of MMD-$k$}]
    $D^{(k)}(\scE_1,\scE_2)=0$ if and only if $\dsE_{\rho\sim\scE_1}[\rho^{\otimes k}] = \dsE_{\sigma\sim\scE_2}[\sigma^{\otimes k}]$; if $D^{(k)}(\scE_1,\scE_2)=0$, $D^{(k')}(\scE_1,\scE_2)=0$ and $\dsE_{\rho\sim\scE_1}[\rho^{\otimes k'}] = \dsE_{\sigma\sim\scE_2}[\sigma^{\otimes k'}]$ for $k \ge k'$. Then, $\scP(\scD^{(k)})\ge\scP(\scD^{(k')})$ for $k\ge k'$. \label{thm: hierarachy of MMD-k}
\end{lemma}


\section*{Retained context: homogeneous polynomials and symmetric tensors}

\begin{lemma}[Homogeneous polynomials and symmetric tensors]
Let $f$ be a homogeneous polynomial of degree $t$ of the elements of $\ket{\psi}\in\mathbb C^d$.
Then there exists a vector $\lvert v_f\rangle \in \mathrm{Sym}^t(\mathbb C^d)$ such that
\begin{equation}
f(\ket{\psi}) = \braket{v_f}{\psi}^{\otimes t}.
\label{eq:app-polynomial-tensor-representation}
\end{equation}
Moreover, for any ensemble $\mathcal E$,
\begin{equation}
\dsE_{\ket{\psi}\sim\mathcal E}\bigl[\lvert f(\ket{\psi})\rvert^2\bigr]
=
\mathrm{Tr}\!\left( \lvert v_f\rangle\langle v_f\rvert \, M_t(\mathcal E)\right).
\label{eq:app-polynomial-moment-identity}
\end{equation}
Hence $M_t(\mathcal E_1)=M_t(\mathcal E_2)$ implies
$\dsE_{\mathcal E_1}[\lvert f\rvert^2]=\dsE_{\mathcal E_2}[\lvert f\rvert^2]$
for all homogeneous degree-$t$ polynomials $f$. \label{lemma: homogeneous and sym}
\end{lemma}


\section*{Retained context: hyperplane separation for distinct pure states}

\begin{lemma}[Hyperplane separation for distinct pure states; $d\ge 2$]
Let $\ket{x},\ket{y}\in\mathbb C^d$ be two distinct unit vectors (pure quantum states), i.e.,
$\ket{x}\neq e^{i\theta}\ket{y}$ for all $\theta\in\mathbb R$. Then there exists a vector $\ket{a}\in\mathbb C^d$
such that
\begin{equation}
\braket{a}{y}=0
\quad\text{and}\quad
\braket{a}{x}\neq 0.
\label{eq:app-hyperplane-separation}
\end{equation}
Equivalently, the linear functional $\ell(\ket{\psi}) := \braket{a}{\psi}$ vanishes at $\ket{y}$ but not at $\ket{x}$. \label{lemma: hyperplane}
\end{lemma}


\section{The counted unit: Theorem A.3 and its complete original proof}

\begin{theorem}[Sufficiency: $k=N$ gives full discriminative power for ensembles with at most $N$ pure states]
Fix $d\ge 2$ and $N\in\mathbb N$.
Let
\[
\mathcal E_1=\{(p_i,\ket{\psi_i})\}_{i=1}^{N_1},
\qquad
\mathcal E_2=\{(q_j,\ket{\phi_j})\}_{j=1}^{N_2}
\]
be two quantum ensembles in $\mathbb C^d$, where $\ket{\psi_i},\ket{\phi_j}$ are unit vectors, all states in each
ensemble are distinct, and $(p_i)_{i=1}^{N_1}$, $(q_j)_{j=1}^{N_2}$ are probability distributions (in particular,
$p_i>0$, $q_j>0$, and $\sum_i p_i=\sum_j q_j=1$). Assume $N_1,N_2\le N$.
If
\[
M_N(\mathcal E_1)=M_N(\mathcal E_2),
\]
then $\mathcal E_1=\mathcal E_2$ up to a permutation of indices (i.e., the two ensembles have the same states with
the same weights).\label{thm: k=N full discriminative power}
\end{theorem}

\begin{proof}
Assume $M_N(\mathcal E_1)=M_N(\mathcal E_2)$. By Lemma~\ref{thm: hierarachy of MMD-k},we have $M_t(\mathcal E_1)=M_t(\mathcal E_2)$ for all integers $1\le t\le N$.

\medskip
\noindent\textit{Step 1: the sets of states coincide.}
Suppose, for contradiction, that there exists a state $\ket{\psi_\star}$ among $\{\ket{\psi_i}\}_{i=1}^{N_1}$
such that $\ket{\psi_\star}\neq e^{i\theta}\ket{\phi_j}$ for every $j\in\{1,\dots,N_2\}$ and every $\theta\in\mathbb R$.

For each $j\in\{1,\dots,N_2\}$, apply Lemma~\ref{lemma: hyperplane}
to the pair $(\ket{x},\ket{y})=(\ket{\psi_\star},\ket{\phi_j})$. This gives a vector $\ket{a_j}\in\mathbb C^d$ such that
\[
\braket{a_j}{\phi_j}=0
\quad\text{and}\quad
\braket{a_j}{\psi_\star}\neq 0.
\]
Define the linear functional $\ell_j(\ket{\psi}) := \braket{a_j}{\psi}$ and the function
\begin{equation}
f(\ket{\psi}) := \prod_{j=1}^{N_2} \ell_j(\ket{\psi})
= \prod_{j=1}^{N_2} \braket{a_j}{\psi}.
\label{eq:app-support-separating-polynomial}
\end{equation}
Each factor $\ell_j(\ket{\psi})$ is a homogeneous polynomial of degree $1$ in the coordinates of $\ket{\psi}$,
so $f$ is a homogeneous polynomial of degree $N_2$.

Now observe:
\begin{itemize}
\item For each $j$, since $\ell_j(\ket{\phi_j})=\braket{a_j}{\phi_j}=0$, we have $f(\ket{\phi_j})=0$.
Therefore,
\[ 
\dsE_{\ket{\phi}\sim\mathcal E_2}\!\big[\Abs{f(\ket{\phi})}^2\big]
= \sum_{j=1}^{N_2} q_j\,\Abs{f(\ket{\phi_j})}^2
=0.
\]
\item On the other hand, $\ell_j(\ket{\psi_\star})=\braket{a_j}{\psi_\star}\neq 0$ for every $j$, hence
$f(\ket{\psi_\star})\neq 0$. Since $p_\star>0$ is the weight of $\ket{\psi_\star}$ in $\mathcal E_1$,
\[
\begin{aligned}
\dsE_{\ket{\psi}\sim\mathcal E_1}\!\big[\Abs{f(\ket{\psi})}^2\big]
&= \sum_{i=1}^{N_1} p_i\,\Abs{f(\ket{\psi_i})}^2\\
&\ge p_\star\,\Abs{f(\ket{\psi_\star})}^2
>0.
\end{aligned}
\]
\end{itemize}

By Lemma~\ref{lemma: homogeneous and sym}, there exists $\ket{v_f}\in\mathrm{Sym}^{N_2}(\mathbb C^d)$
such that
\[
\begin{gathered}
\dsE_{\ket{\psi}\sim\mathcal E}\!\big[\Abs{f(\ket{\psi})}^2\big]
=
\mathrm{Tr}\!\Big(\ketbra{v_f}{v_f}\,M_{N_2}(\mathcal E)\Big)
\\
\text{for any ensemble }\mathcal E.
\end{gathered}
\]
Hence the strict inequality
\[
\dsE_{\mathcal E_1}\!\big[\Abs{f}^2\big] \neq \dsE_{\mathcal E_2}\!\big[\Abs{f}^2\big]
\]
implies
\[
\mathrm{Tr}\!\Big(\ketbra{v_f}{v_f}\,M_{N_2}(\mathcal E_1)\Big)
\neq
\mathrm{Tr}\!\Big(\ketbra{v_f}{v_f}\,M_{N_2}(\mathcal E_2)\Big),
\]
and therefore $M_{N_2}(\mathcal E_1)\neq M_{N_2}(\mathcal E_2)$.
But $N_2\le N$ and we already have $M_t(\mathcal E_1)=M_t(\mathcal E_2)$ for all $t\le N$, a contradiction.
Thus, every state in $\mathcal E_1$ must coincide with some state in $\mathcal E_2$ up to global phase.

By symmetry (swapping the roles of $\mathcal E_1$ and $\mathcal E_2$), every state in $\mathcal E_2$ must also coincide
with some state in $\mathcal E_1$ up to global phase. Therefore the two ensembles contain the same set of states (up to
a permutation of indices and global phases), and in particular $N_1=N_2$.

\medskip
\noindent\textit{Step 2: the weights coincide.}
After relabeling, assume $\ket{\psi_i}=e^{i\theta_i}\ket{\phi_i}$ for all $i\in\{1,\dots,N_1\}$.
Fix an index $i$. For each $j\neq i$, apply Lemma~\ref{lemma: hyperplane}
to the pair $(\ket{x},\ket{y})=(\ket{\psi_i},\ket{\psi_j})$ to obtain $\ket{a_{ij}}$ such that
\[
\braket{a_{ij}}{\psi_j}=0
\quad\text{and}\quad
\braket{a_{ij}}{\psi_i}\neq 0.
\]
Define $\ell_{ij}(\ket{\psi}) := \braket{a_{ij}}{\psi}$ and
\begin{equation}
f_i(\ket{\psi})
:=\prod_{j\neq i}\frac{\ell_{ij}(\ket{\psi})}{\ell_{ij}(\ket{\psi_i})}
=\prod_{j\neq i}\frac{\braket{a_{ij}}{\psi}}{\braket{a_{ij}}{\psi_i}}.
\label{eq:app-weight-interpolation-polynomial}
\end{equation}
Then $f_i(\ket{\psi_i})=1$ and $f_i(\ket{\psi_j})=0$ for all $j\neq i$.
Consequently,
\[
\dsE_{\ket{\psi}\sim\mathcal E_1}\!\big[\Abs{f_i(\ket{\psi})}^2\big]
= \sum_{r=1}^{N_1} p_r\,\Abs{f_i(\ket{\psi_r})}^2
= p_i,
\]
and similarly,
\[
\dsE_{\ket{\phi}\sim\mathcal E_2}\!\big[\Abs{f_i(\ket{\phi})}^2\big]
= q_i,
\]
since $\Abs{f_i(e^{i\theta}\ket{\psi})}=\Abs{f_i(\ket{\psi})}$ and the two ensembles share the same states up to phases.

Because $f_i$ is a homogeneous polynomial of degree $N_1-1$ divided by a nonzero constant,
it is (as a function of coordinates) homogeneous of degree $N_1-1\le N-1$.
Applying Lemma~\ref{lemma: homogeneous and sym} to $f_i$ gives
\[
\dsE_{\mathcal E}\!\big[\Abs{f_i}^2\big]
=
\mathrm{Tr}\!\Big(\ketbra{v_{f_i}}{v_{f_i}}\,M_{N_1-1}(\mathcal E)\Big).
\]
Since $N_1-1\le N$ and $M_{N_1-1}(\mathcal E_1)=M_{N_1-1}(\mathcal E_2)$, we conclude
\begin{equation}
p_i=\dsE_{\mathcal E_1}\!\big[\Abs{f_i}^2\big]
=\dsE_{\mathcal E_2}\!\big[\Abs{f_i}^2\big]
=q_i.
\label{eq:app-ensemble-weight-equality}
\end{equation}
As $i$ was arbitrary, $p_i=q_i$ for all $i$, proving that the two ensembles are identical up to permutation.
\end{proof}

\end{document}
